\section{Interpretacja uzyskanych wyników}
\label{sec:InterpretacjaWynikow}

Uzyskane wyniki pokazują, że różnice wydajnościowe między algorytmami wynikały zarówno z~ich charakterystyki obliczeniowej,
jak i~z~przyjętego sposobu implementacji obliczeń równoległych.

Na obserwowane rezultaty wpływały między innymi sposób podziału danych, liczba tworzonych procesów, 
moment przejścia do wykonania sekwencyjnego, a~także dodatkowy narzut związany z~tworzeniem i~synchronizacją procesów.
Wpływ tych czynników różnił się pomiędzy poszczególnymi algorytmami, co było widoczne w uzyskanych czasach wykonania, 
wartościach przyspieszenia i~efektywności oraz wykorzystaniu zasobów.

\subsection*{Wpływ sposobu organizacji obliczeń równoległych}

Istotnym czynnikiem wpływającym na uzyskane wyniki był sposób sterowania poziomem równoległości. 
W przypadku Quick Sort oraz Merge Sort maksymalną głębokość równoległego podziału określano zgodnie z~zależnością

\begin{equation}
    \texttt{max\_depth} = \left\lfloor \log_2(p) \right\rfloor.
\end{equation}

gdzie $p$ oznacza liczbę jednostek wykonawczych. 

Dodatkowo wprowadzono progi minimalnego rozmiaru fragmentu danych, po których osiągnięciu dalsze przetwarzanie realizowano sekwencyjnie. 
W~przypadku osiągnięcia maksymalnej głębokości również następowało przejście do wykonania sekwencyjnego. 
Oznacza to, że utworzenie kolejnego procesu było możliwe tylko w~przypadku niespełnienia obu tych warunków.

W~Bucket Sort oraz Sample Sort zastosowano natomiast progi określające uruchomienie 
wykonania równoległego oraz utworzenie procesu dla określonej grupy danych. 
Rozwiązania te ograniczały narzut związany z~uruchamianiem procesów dla niewielkich zadań, 
jednocześnie wpływając na rzeczywisty poziom równoległości.

Przyjęte wartości progów zależały od konfiguracji liczby jednostek wykonawczych,
w~związku z~czym nie zapewniały jednakowo korzystnego podziału pracy dla każdego rozmiaru i~rodzaju danych.
W~części analizowanych konfiguracji obserwowano jednocześnie niższe wykorzystanie CPU, mniejsze przyspieszenie oraz niższą efektywność.

Dodatkowa diagnostyka pozwoliła określić rzeczywistą liczbę procesów potomnych tworzonych podczas wykonania algorytmów. 
Wyniki te pokazują, że zwiększanie liczby zadeklarowanych jednostek wykonawczych
nie prowadziło do proporcjonalnego zwiększenia liczby tworzonych procesów.

W przypadku Quick Sort partycjonowanie prowadziło do powstawania fragmentów o~różnych rozmiarach, 
a~tym samym do nierównomiernego rozłożenia pracy pomiędzy procesami. 
Dodatkowo możliwość tworzenia kolejnych procesów ograniczały maksymalna głębokość podziału oraz przyjęty próg minimalnego rozmiaru fragmentu.

Potwierdza to diagnostyka, w~której po zwiększeniu liczby jednostek wykonawczych z~8 do 16 i~32 
liczba utworzonych procesów dla Quick Sort pozostała na poziomie 4. 
Dla zbioru \texttt{random\_int} najlepszy czas uzyskano przy 2 jednostkach wykonawczych,
natomiast przy 8, 16 i~32 jednostkach czas wykonania nie ulegał dalszemu skróceniu.

W przypadku Merge Sort podział danych miał bardziej regularny charakter, ponieważ na kolejnych 
poziomach przetwarzany zakres dzielono na dwie części. 
Po zakończeniu sortowania fragmentów konieczne było jednak oczekiwanie na procesy potomne oraz scalanie uzyskanych części. 
Wraz ze zwiększaniem poziomu równoległości rosło więc znaczenie kosztów tworzenia i~synchronizacji procesów oraz operacji scalania.

Diagnostyka wykazała, że przy 8, 16 i~32 jednostkach wykonawczych dla Merge Sort utworzono po 2 procesy potomne. 
Jest to zgodne z~ograniczeniem wynikającym z~maksymalnej głębokości oraz progów przejścia do wykonania sekwencyjnego. 
Najlepszy czas uzyskano przy 4 jednostkach wykonawczych,
natomiast przy 8, 16 i~32 jednostkach liczba utworzonych procesów pozostawała na poziomie 2,
a~dalsze zwiększanie ich liczby nie prowadziło już do poprawy wyniku.

W przypadku Bucket Sort dane przydzielano do kubełków, które następnie grupowano i~przekazywano dostępnym procesom. 
Przyjęte progi decydowały o~wykorzystaniu wykonania równoległego oraz utworzeniu oddzielnego procesu dla danej grupy. 
Pozwalało to na przetwarzanie mniejszych zadań bez ponoszenia dodatkowych kosztów związanych z~wieloprocesowością.

Diagnostyka wykazała, że w przypadku Bucket Sort liczba utworzonych procesów była zgodna z~liczbą 
zadeklarowanych jednostek wykonawczych i~wynosiła odpowiednio 2, 4, 8, 16 oraz 32. 
Wraz ze zwiększaniem liczby procesów obserwowano skracanie czasu wykonania aż do konfiguracji z~16 jednostkami wykonawczymi. 
Dalsze zwiększenie ich liczby do 32 nie przyniosło już poprawy, mimo utrzymującego się wysokiego wykorzystania procesora. 

Sample Sort wykorzystywał wieloetapową organizację obliczeń, obejmującą równoległe sortowanie fragmentów, 
wyznaczanie elementów podziałowych, redystrybucję danych oraz sortowanie utworzonych grup. 
Wiązało się to z~kilkoma etapami tworzenia i~synchronizacji procesów oraz dodatkową pracą związaną z~podziałem danych.

W~diagnostyce Sample Sort tworzył procesy na dwóch etapach przetwarzania, obejmujących sortowanie lokalne oraz sortowanie grup kubełków.
W~każdej z~faz uruchamiano liczbę procesów odpowiadającą zadeklarowanej liczbie jednostek wykonawczych, 
dlatego łączna liczba utworzonych procesów wynosiła odpowiednio 4, 8, 16, 32 i~64 dla 2, 4, 8, 16 i~32 jednostek. 
Procesy z~poszczególnych faz nie były jednak uruchamiane jednocześnie.

Najlepszy wynik uzyskano przy 8 jednostkach wykonawczych, natomiast dalsze zwiększanie ich liczby do 16 i~32 nie prowadziło już do poprawy czasu wykonania.
Przy większej liczbie jednostek ulegała zmniejszeniu ilość danych przypadająca na pojedyncze zadanie, 
podczas gdy koszty organizacji kolejnych etapów pozostawały obecne.

\subsection*{Wpływ narzutu wieloprocesowości}

Na uzyskane wyniki wpływał również narzut związany z~wykorzystaniem wieloprocesowości. 
Badane warianty równoległe zaimplementowano z~wykorzystaniem modułu \texttt{multiprocessing}, tworząc procesy klasą \texttt{Process}. 
Współdzielenie danych realizowano poprzez mechanizm \texttt{shared\_memory}.

Tworzenie procesów i~kontrolowanie ich wykonania wiązało się z~dodatkowym narzutem, który nie występował w~wariantach sekwencyjnych. 
Obejmował on między innymi uruchamianie procesów, oczekiwanie na ich zakończenie oraz operacje związane z~obsługą pamięci współdzielonej.
Ponadto w~części implementacji konieczne było kopiowanie fragmentów danych do struktur 
lokalnych na potrzeby dalszego przetwarzania, co stanowiło dodatkowy element wykonywanych operacji.

Znaczenie tego narzutu rosło wraz ze zmniejszaniem ilości pracy przypadającej na pojedynczy proces. 
Niewielkie korzyści z~równoległości obserwowano dla małych zbiorów danych,
natomiast wraz ze zwiększaniem liczby jednostek wykonawczych obserwowano spadek efektywności,
przy jednoczesnym zmniejszaniu ilości pracy przypadającej na pojedynczy proces.

Dodatkowe znaczenie miało środowisko przeprowadzania eksperymentów. Badania wykonano w~systemie Windows, 
gdzie moduł \texttt{multiprocessing} wykorzystuje metodę \texttt{spawn} do uruchamiania nowych procesów. 
Metoda ta prowadzi do~uruchomienia odrębnego procesu interpretera Pythona, 
co wiąże się z~dodatkowym narzutem inicjalizacji przed rozpoczęciem właściwych obliczeń.

W~opracowanych implementacjach nie stosowano stałej puli procesów.
Procesy tworzono w~momentach, w~których poszczególne etapy algorytmów wymagały wykonania zadań równoległych. 
Koszty ich uruchamiania i~synchronizacji stanowiły dodatkowy element czasu wykonania,
szczególnie przy niewielkiej ilości pracy przypadającej na pojedynczy proces.

Dodatkowych informacji na temat przebiegu obliczeń dostarczyło profilowanie przeprowadzone z~wykorzystaniem narzędzia \texttt{py-spy}.
Na wygenerowanych wykresach płomieni widoczne były między innymi stosy wywołań związane z~mechanizmem \texttt{multiprocessing}, 
w tym funkcje \texttt{spawn\_main} oraz \texttt{\_bootstrap}, a także funkcje odpowiadające właściwym operacjom analizowanych algorytmów.
Przykładowy wykres dla równoległej implementacji Merge Sort przedstawiono na~rysunku~\ref{fig:merge-sort-flamegraph}.

\begin{figure}[H]
    \centering
    \includegraphics[width=\textwidth]
    {images/chapter7/Merge_Sort_-_ParallelV4.png}
    \caption{Wykres płomieni dla równoległej implementacji Merge Sort uzyskany za pomocą narzędzia \texttt{py-spy}}
    \label{fig:merge-sort-flamegraph}
\end{figure}

Obecność funkcji związanych z~uruchamianiem procesów potomnych na wykresie odpowiada realizacji obliczeń z~użyciem modułu \texttt{multiprocessing}.
Samego udziału funkcji \texttt{spawn\_main} nie należy jednak wiązać bezpośrednio z~czasem przeznaczonym wyłącznie na utworzenie procesu, 
ponieważ jest ona częścią stosu wywołań procesu potomnego, w~obrębie którego realizowane są również kolejne operacje algorytmu.

Uzyskane wyniki należy zatem interpretować jako rezultat zarówno właściwości poszczególnych algorytmów, 
jak i~przyjętego sposobu realizacji obliczeń równoległych. 
Istotne znaczenie miały sposób podziału danych, wielkość wykonywanych zadań, 
zastosowane progi przejścia do wariantu sekwencyjnego oraz narzut związany z~tworzeniem i~synchronizacją procesów.